\documentclass[11pt]{article}
\usepackage[margin=1in]{geometry}
\usepackage{hyperref}
\usepackage{microtype}
\hypersetup{colorlinks=true,linkcolor=blue,citecolor=blue,urlcolor=blue}
\title{A Structured DDJ-FLX4 Curriculum for Chinese-Speaking Beginner DJs\\\large Curriculum architecture, practice progression, and evidence boundaries}
\author{Stan Chen}
\date{23 November 2026\\[4pt]\small Unpublished working draft; not submitted to arXiv}
\begin{document}
\maketitle
\begin{center}\small This draft has no arXiv identifier. Claims are limited to the evidence described below.\end{center}
\begin{abstract}
This design report describes a Chinese-language beginner curriculum organized around the Pioneer DJ DDJ-FLX4 and rekordbox. The published sequence contains fourteen technical lessons and five career-oriented lessons, progressing from setup and cue monitoring through beatmatching, transitions, set preparation, performance contexts, and early professional practice. The curriculum connects procedural equipment knowledge with listening and practice. This is an instructional design description, not an educational efficacy study: no controlled comparison, learner outcome dataset, retention assessment, or representative accessibility evaluation is reported. We define a learner-centered evaluation plan and distinguish transferable DJ skills from controller-specific instruction.
\end{abstract}
\section{Motivation and audience}
The curriculum offers a connected learning path for Chinese-speaking beginners who own or are considering a DDJ-FLX4 \cite{r1}. It uses rekordbox \cite{r2} for concrete examples while treating listening, phrasing, selection, preparation, and performance judgment as skills that can transfer to other setups.\par
A complete sequence of lessons is an instructional artifact, not evidence that all learners progress at the same rate or achieve the intended skills.\par
\section{Curriculum structure}
The course contains nineteen lessons \cite{r3}: fourteen technical lessons and five career-oriented lessons. The technical sequence covers setup, audio routing, gain, EQ, headphone monitoring, manual beatmatching, phrase structure, transitions, loops, hot cues, effects, harmonic mixing, livestream setup, practice, and set construction. The final lessons address early performance contexts and preparation.\par
The progression moves from a controllable practice environment toward contextual decisions. Learners establish the signal chain, practice timing and transitions, and then assemble and perform a set.\par
\section{Limitations and future work}
There is no learner cohort, controlled teaching comparison, pre/post assessment, or verified long-term retention study in this draft. Career advice is contextual and is not a guarantee of paid work.\par
A first evaluation could assess a recorded transition before and after the course with blinded rubric ratings for beat alignment, phrase timing, level management, and continuity. Report participant experience, practice exposure, attrition, and retention.\par
\section{Conclusion}
The curriculum provides a structured starting point for DDJ-FLX4 beginners. Its educational value remains an empirical question for learner-centered evaluation.\par
\begin{thebibliography}{99}
\bibitem{r1} AlphaTheta Corporation. DDJ-FLX4 official product information. \url{https://www.pioneerdj.com/en/product/controller/ddj-flx4/black/overview/}
\bibitem{r2} AlphaTheta Corporation. rekordbox official product information. \url{https://rekordbox.com/en/}
\bibitem{r3} Chen, S. DJ Notes repository. \url{https://github.com/stancsz/djnotes}
\end{thebibliography}
\end{document}
